\label{summary}

Human genetic studies and analyses of diseased tissue provide different views of metabolic dysfunction-associated steatotic liver disease (MASLD), but connecting these observations to experiments remains challenging. Here, we assemble a human MASLD resource that records liver-trait genetics beside disease-associated liver expression, with cell and tissue context. Most colocalized genetic signals in the resource come from liver enzymes such as serum alanine aminotransferase (ALT), the common clinical marker of liver injury (417 of 430 primary multi-signal genes after allele-aware, positive-semidefinite LD-projected eQTL refitting). In a genetic covariance model, 0.701 (95\% interval 0.601 to 0.775; replication 0.802) of ALT heritability is shared with neither liver fat nor cirrhosis. Of 286 genes with an allele-direction call, 260 have liver-enzyme support only (234 of 253 in the non-circular European enzyme-defined subset), although most are underpowered for liver fat. Across 167 colocalized ALT/AST loci, the enzyme-raising allele raises plasma levels of other hepatocyte proteins in aggregate (IVW slope 0.333 protein SD per MVP ALT SD with zero sample overlap; 0.435 per Pan-UKBB ALT SD; 26 of 167 loci pass individually). At TNFSF10, the ALT signal colocalizes with plasma TRAIL (SuSiE PP.H4 1.00; 4-trait exact ABF PP 0.971; zero coding proxies at r$^2 \ge$ 0.8). The few ALT loci that are also associated with liver fat colocalize with plasma lipid species more often than other ALT loci (4 of 6 against 10 of 143; odds ratio 25.2, \emph{p} = 6.8 $\times$ 10$^{-4}$; 3 of 5 against 11 of 144 for steatosis-concordant loci, \emph{p} = 0.006; attenuated to odds ratio 5.45, \emph{p} = 0.053 at PP.H4 $\ge$ 0.5). In 28 drug mechanisms tested in randomized MASLD or MASH trials, no comparison of target-gene genetics with the matched trial endpoint was informative. A sequence model (AlphaGenome 2, evaluated on the ALL\_FOLDS checkpoint whose pretraining may overlap liver chromatin tracks) predicts measured liver chromatin-accessibility QTL effects (Spearman 0.719), and we use it only to annotate accessibility at fine-mapped variants. In tissue, multicohort transcriptomics describes cross-sectional stage-associated remodeling; of 7,214 stage genes replicated in 493 held-out livers, 1,276 change inside hepatocytes in uncorrected single-nucleus pseudobulk (188 after empirical ambient-RNA subtraction, 0 of 9 negative controls called), and GNMT, MAT1A and CYP2C19 are lower per hepatocyte in one held-out single-nucleus cohort whose controls were collected separately while two smaller held-out cohorts were uninformative. We organize these observations in the MASLD Gene Catalog, which records the evidence available for each gene, identifies missing measurements and proposes experiments to resolve the remaining questions.
